\section{The family \texorpdfstring{$\{0,a,b,2b\}$}{\{0,a,b,2b\}}}\label{sec:family}

This section proves Theorem~\ref{thm:C}. At the modulus $b$ the pairs of the family are similarly ordered, its missing positions are complements of arithmetic progressions on $\Z/b\Z$, and the spectrum becomes a sum over circular gaps, which the three-distance theorem evaluates.

Throughout this section $0<a<b$ are coprime and $A_0=\{0,a,b,2b\}$, so $d=2b$, $S_1=\langle a,b\rangle$ and $S_2=\langle b,2b-a\rangle$. We put $\M_4(a,b)=\M(a,b)+\M(b,2b-a)$ and $\phi(j)=\phi_{b,a}(j)$ (Section~\ref{sec:prelim}), so $\phi(1)=\binom b2$ and $\phi(j)=0$ for $j\ge b$.

\subsection{Structure and spectrum}\label{subsec:family-structure}

\begin{lemma}\label{lem:family-ci}
Over every field, the ideal of $C_A$ is generated by $x_2^2-x_0x_3$ and $x_1^b-x_2^{\,j}x_0^{\,z}x_3^{\,w}$, where $x_0,\dots,x_3$ correspond to the parts $0,a,b,2b$, $j=a\bmod2$, $w=(a-j)/2$ and $z=(2b-a-j)/2$. So $C_A$ is a complete intersection of degrees $2$ and $b$ and $A_0$ is CM; $|hA_0|=(h+1)^2$ for $h\le b-1$ and $|hA_0|=2bh+1-(b-1)^2$ for $h\ge b-1$; $c_1=(a-1)(b-1)$, $c_2=(b-1)(2b-a-1)$ and $\sigma(A)=b-1$.
\end{lemma}

\begin{proof}
Both binomials are homogeneous (of degrees $2$ and $b=j+w+z$) and lie in the ideal $I$ of $C_A$. Modulo them every monomial is congruent to a monomial $x_0^rx_1^ex_2^\varepsilon x_3^t$ of the same degree with $e<b$, $\varepsilon\le1$: replacing $x_1^b$ lowers the exponent of $x_1$, replacing $x_2^2$ keeps it and lowers that of $x_2$. Such a monomial maps to the point with first coordinate $ae+b\varepsilon+2bt$; for two of them of equal degree and equal image, $ae\equiv ae'\pmod b$ gives $e=e'$, then parity gives $\varepsilon=\varepsilon'$ and $t=t'$, and the degree gives $r=r'$. So in each degree these monomials have linearly independent images, every homogeneous element of $I$ reduces to $0$, and $I$ is generated by the two binomials; a prime of height two with two generators is a complete intersection (compare \cite{BermejoGarciaMarco2014}). Hence $|hA_0|$ is the number of these monomials of degree $h$, namely $\sum_{e=0}^{\min(h,b-1)}(2(h-e)+1)$. The conductors are those of $\langle a,b\rangle$ and $\langle b,2b-a\rangle$, $c_1+c_2=2(b-1)^2$, and $\sigma(A)=\lceil(b-1)^2/b\rceil=b-1$ (Section~\ref{sec:prelim}).
\end{proof}

The integer $b$ is a modulus (Section~\ref{subsec:counting}), with $N=2h$.

\begin{lemma}[Residue intervals modulo $b$]\label{lem:family-residue}
For $0\le s<b$ let $i_s\in[0,b-1]$ satisfy $i_sa\equiv-s\pmod b$. At $m=b$, $u_s=(s+i_sa)/b=\lceil i_sa/b\rceil$ and $v_s=2i_s-u_s$, so $kb-s\in hA_0$ if and only if $u_s\le k\le2h-v_s$. As functions of $i=i_s$, $u$ is nondecreasing with steps $0,1$ and $v$ is increasing with steps $1,2$. Hence $(u,v)$ is similarly ordered at $m=b$, $U=a$, $V=2b-a-2$, and $\Tot$ is nondecreasing and concave and equals $\M_4(a,b)$ exactly for $h\ge b-1$.
\end{lemma}

\begin{proof}
The Ap\'ery elements modulo $b$ are $ia$ for $S_1$ and $i(2b-a)\equiv-ia$ for $S_2$, so $w^{(1)}_{-s}=i_sa$ and $w^{(2)}_s=i_s(2b-a)$, which gives $u_s$ and $v_s$; Lemma~\ref{lem:residue} applies by Lemma~\ref{lem:family-ci}. Since $0<a/b<1$, $\lceil(i+1)a/b\rceil-\lceil ia/b\rceil\in\{0,1\}$. If $u_s>u_t$ then $i_s>i_t$, because $u$ is nondecreasing in $i$, and then $v_s>v_t$, because $v$ is increasing. Both maxima occur at $i=b-1$. By Theorem~\ref{thm:spectrum}(3), the mass is reached exactly when $2h\ge U+V-1=2b-3$.
\end{proof}

The family need not be similarly ordered at $d=2b$ (Section~\ref{subsec:similar}); at the modulus $b$, Proposition~\ref{prop:threshold} gives the level $\sigma(A)=b-1$ of Lemma~\ref{lem:family-ci} again, and $\reg\Bbbk[\mathcal S]=b$, the regularity of a complete intersection of degrees $2$ and $b$.

For $1\le i\le b-1$ let $s(i)\in[1,b-1]$ satisfy $s(i)\equiv-ia\pmod b$, put $\mathcal P_j=\{s(i):j\le i\le b-1\}$, and let $\mathcal R_j$ be the multiset of the numbers $g-1\ge1$ over the gaps $g$ of $\{0,a,\dots,(j-1)a\}$ on $\Z/b\Z$.

\begin{theorem}[Run lengths]\label{thm:family-runs}
For $h\ge1$ the lengths of the maximal runs of integers of $[0,2hb]$ outside $hA_0$ form the multiset
\[
\mathcal R_1\sqcup\bigsqcup_{j=2}^{h}2\mathcal R_j\sqcup\mathcal R_{h+1}\quad(h<b-1),\qquad \mathcal R_1\sqcup\bigsqcup_{j=2}^{b-1}2\mathcal R_j\quad(h\ge b-1),
\]
where $2\mathcal R_j$ is two copies. In particular the distribution of $D_h$ on $[0,2hb]$ is unchanged when $a$ is replaced by $b-a$.
\end{theorem}

\begin{proof}
By Lemma~\ref{lem:family-residue}, position $s(i)$ is present in block $k$ exactly for $k\in E_i=[\alpha_i,\alpha_i+2h-2i]$, $\alpha_i=\lceil ia/b\rceil$, an interval inside $[1,2h]$ when nonempty. As $\alpha_{i+1}-\alpha_i\in\{0,1\}$, $E_{i+1}\subseteq E_i$, and $|E_i|=\max(2(h-i)+1,0)$. So the holes of block $k$ are $\mathcal P_j$ with $j=1+\max(\{0\}\cup\{i:k\in E_i\})$. Since $E_1=[1,2h-1]$, one block has holes $\mathcal P_1$. For $2\le j\le\min(h,b-1)$ the two blocks of $E_{j-1}\setminus E_j$ have holes $\mathcal P_j$. If $h<b-1$, the block of $E_h$ has holes $\mathcal P_{h+1}$, and if $h\ge b-1$, the blocks of $E_{b-1}$ have none. The complement of $\mathcal P_j$ in $[1,b-1]$, together with $0$, is $\{0,-a,\dots,-(j-1)a\}$, and a gap $g$ of this set is a run of $g-1$ holes. The map $x\mapsto-x$ carries the gaps of $\{0,a,\dots,(j-1)a\}$ to those of this set, so the runs of $\mathcal P_j$ have lengths $\mathcal R_j$. Replacing $a$ by $b-a$ reflects each point set and preserves each $\mathcal R_j$, and a run of length $L$ carries the undershoots $1,\dots,L$.
\end{proof}

\begin{theorem}[Spectrum of the family]\label{thm:family-spectrum}
At $m=b$ the spectrum of $A_0$ is $\binom b2$ once and $\phi(j)$ twice for $2\le j\le b-1$, and $\phi$ is nonincreasing. Hence
\begin{equation}\label{eq:family-tot}
\Tot(h)=\binom b2+2\sum_{j=2}^{h}\phi(j)+\phi(h+1)\qquad(1\le h\le b-1),
\end{equation}
and $\Tot(h)=\M_4(a,b)$ for $h\ge b-1$.
\end{theorem}

\begin{proof}
Write $u(i)=\lceil ia/b\rceil$ and $v(i)=2i-u(i)$ for the pairs of Lemma~\ref{lem:family-residue} as functions of $i$. Then $B_k=\mathcal P_j$ for the $u(j)-u(j-1)$ values $u(j-1)\le k<u(j)$, and $R_\ell=\mathcal P_j$ for the $v(j)-v(j-1)$ values $v(j-1)\le\ell<v(j)$. The two counts add up to $2$ for $j\ge2$, while $\mathcal P_1$ occurs only as $R_0$, since $k\ge1=u(1)$. By the proof of Theorem~\ref{thm:family-runs}, $f(\mathcal P_j)=\phi(j)$. Adding a point to a circular set splits a gap $x+y$ and lowers $\binom{x+y}2$ by $xy\ge0$, so $\phi$ is nonincreasing, and Theorem~\ref{thm:spectrum}(3) gives $\Tot(h)$ as the sum of the $2h$ largest entries.
\end{proof}

\begin{corollary}[Mass of the family]\label{cor:family-mass}
For coprime $0<a<b$, $\M_4(a,b)=2\Phi(b,a)-\binom b2$.
\end{corollary}

\begin{proof}
By Lemma~\ref{lem:spectrum-mass} the mass is the sum of the spectrum, $\binom b2+2\sum_{j=2}^{b-1}\phi(j)=2\Phi(b,a)-\binom b2$, as $\phi(1)=\binom b2$ and $\phi(b)=0$.
\end{proof}

\subsection{The whole range through continued fractions}\label{subsec:family-cf}

Let $a/b=[0;\theta_1,\dots,\theta_K]$ with $\theta_K\ge2$, with convergents $p_j/q_j$, $p_{-1}=1$, $q_{-1}=0$, $p_0=0$, $q_0=1$, so $q_K=b$, and put $\varepsilon_j=|q_ja-p_jb|$ for $-1\le j\le K$. Then $\varepsilon_{-1}=b$, $\varepsilon_0=a$, $\varepsilon_{K-1}=1$, $\varepsilon_K=0$ and $\varepsilon_{j-1}=\theta_{j+1}\varepsilon_j+\varepsilon_{j+1}$. The intervals $[\nu q_k+q_{k-1},(\nu +1)q_k+q_{k-1})$, $1\le \nu \le \theta_{k+1}$, tile $[q_k+q_{k-1},q_{k+1}+q_k)$, and these tile $[1,b+q_{K-1})$ for $0\le k<K$. So each $h\in[2,b]$ has a unique expression
\begin{equation}\label{eq:cf-index}
h-1=\nu q_k+q_{k-1}+r,\qquad 0\le k\le K-1,\ 1\le \nu \le \theta_{k+1},\ 0\le r<q_k ,
\end{equation}
and the block $(k,\nu )=(K-1,\theta_K)$, which would start at $h=b+1$, does not occur. The $h$ with a given $(k,\nu )$ form a \emph{block}.

The next theorem is the rational form of the classical three-distance theorem, in the notation above.

\begin{theorem}[Three distances]\label{thm:three-distance}
For $2\le h\le b$ as in \eqref{eq:cf-index}, the gaps of $\{0,a,\dots,(h-1)a\}$ on $\Z/b\Z$ are $h-q_k$ gaps $\varepsilon_k$, $r+1$ gaps $\varepsilon_{k-1}-\nu \varepsilon_k$ and $q_k-r-1$ gaps $\varepsilon_{k-1}-(\nu -1)\varepsilon_k$. Hence
\[
\phi(h)=(h-q_k)\binom{\varepsilon_k}2+(r+1)\binom{\varepsilon_{k-1}-\nu \varepsilon_k}2+(q_k-r-1)\binom{\varepsilon_{k-1}-(\nu -1)\varepsilon_k}2 .
\]
\end{theorem}

\begin{proof}
For irrational $\alpha$ with convergents $p_j/q_j$ and $\eta_j=|q_j\alpha-p_j|$, the three distance theorem in the form of \cite[Sect.~3]{AlessandriBerthe1998} (due to S\'os \cite{Sos1958}, Sur\'anyi \cite{Suranyi1958} and \'Swierczkowski \cite{Swierczkowski1958}) states that for $n=\nu q_k+q_{k-1}+r$ as above the points $0,\alpha,\dots,n\alpha$ cut $\mathbb R/\Z$ into $n+1-q_k$ arcs $\eta_k$, $r+1$ arcs $\eta_{k-1}-\nu \eta_k$ and $q_k-r-1$ arcs $\eta_{k-1}-(\nu -1)\eta_k$. Let $\alpha_N=[0;\theta_1,\dots,\theta_K,N,1,1,\dots]$: its convergents of index at most $K$ are those of $a/b$, and $\alpha_N\to a/b$. Take $n=h-1$; its expression involves only $k\le K-1$ and $\nu \le \theta_{k+1}$, so it is the same for $\alpha_N$. The points $ia/b\bmod1$, $0\le i<h\le b$, are distinct with circular spacing at least $1/b$. Once $|\alpha_N-a/b|<1/(3b^2)$, each $i\alpha_N$ lies within $1/(3b)$ of $ia/b$, so the circular order is the same and each arc converges to the corresponding rational gap divided by $b$, while $\eta_j(\alpha_N)\to\varepsilon_j/b$ for $j\le K-1$. Pass to the limit and multiply by $b$.
\end{proof}

Thus the formula for $\phi$ changes only from one block to the next. On a block, raising $h$ splits a gap $x+\varepsilon_k$, $x=\varepsilon_{k-1}-\nu \varepsilon_k$, into $x$ and $\varepsilon_k$, so $\phi$ is affine there with slope $-\delta_{k,\nu }$, $\delta_{k,\nu }=\varepsilon_k(\varepsilon_{k-1}-\nu \varepsilon_k)$. The quadratic pieces of $\Tot$ are shifted.

\begin{corollary}[Quadratic pieces]\label{cor:quadratic-pieces}
Put $\Tot(0)=0$, and let $h^*=\nu q_k+q_{k-1}+1$ be the first $h$ of block $(k,\nu )$. Then
\[
\Tot(h^*-2+j)=\Tot(h^*-2)+2\phi(h^*-1)\,j-\delta_{k,\nu }\,j^2\qquad(0\le j\le q_k),
\]
and the intervals $[h^*-2,h^*+q_k-2]$ tile $[0,b-1]$ with shared endpoints.
\end{corollary}

\begin{proof}
The formula of Theorem~\ref{thm:three-distance} at $r=-1$ equals $\phi(h^*-1)$. For $\nu >1$ it is the value of block $(k,\nu -1)$ at $r=q_k-1$. For $\nu =1$ and $k\ge1$ it is $q_{k-1}\binom{\varepsilon_k}2+q_k\binom{\varepsilon_{k-1}}2$, the value at the last $h=q_k+q_{k-1}$ of block $(k-1,\theta_k)$, because $\varepsilon_{k-2}-\theta_k\varepsilon_{k-1}=\varepsilon_k$. For $k=0$ and $\nu =1$ both are $\binom b2$. So $\phi(h^*-1+i)=\phi(h^*-1)-\delta_{k,\nu }i$ for $0\le i\le q_k$. By \eqref{eq:family-tot}, $\Tot(h)-\Tot(h-1)=\phi(h)+\phi(h+1)$ for $1\le h\le b-1$, which at $h=h^*-2+j$, $1\le j\le q_k$, is $2\phi(h^*-1)-\delta_{k,\nu }(2j-1)$, and summing over $j$ gives the formula. Consecutive blocks start at $h^*$ and $h^*+q_k$, also across rows, as $\theta_{k+1}q_k+q_{k-1}+q_k=q_{k+1}+q_k$. The first block has $h^*=2$ and the last ends at $h=b$.
\end{proof}

For $a=2$, $b=9$ ($[0;4,2]$) the last block of $\phi$ is $h=6,\dots,9$, but $\Tot(6),\dots,\Tot(9)=132,135,136,136$ has third difference $1$; the last quadratic piece of $\Tot$ is $4\le h\le8$, with values $120,127,132,135,136$.
